\documentclass[10pt,a4paper]{article}
\usepackage[T1]{fontenc}
\usepackage[utf8]{inputenc}
\usepackage{lmodern,amsmath,amssymb,booktabs,array,microtype}
\pdfmapfile{+lm.map}\pdfmapfile{+cm.map}\pdfmapfile{+cmextra.map}\pdfmapfile{+symbols.map}
\pdfinfoomitdate=1\pdftrailerid{}
\usepackage[a4paper,margin=23mm]{geometry}
\usepackage{xcolor,hyperref,fancyhdr}
\definecolor{ink}{HTML}{183A4A}
\hypersetup{colorlinks=true,linkcolor=ink,urlcolor=ink,pdftitle={CST-DUST: Erweiterte Verklebungsregel v0.1}}
\pagestyle{fancy}\fancyhf{}\fancyhead[L]{\small CST-DUST | Erweiterter Randvertrag}
\fancyhead[R]{\small v0.1 | 11.11.2026}\fancyfoot[C]{\thepage}
\setlength{\headheight}{14pt}\setlength{\parindent}{0pt}\setlength{\parskip}{5pt}
\setlength{\emergencystretch}{2em}
\newcommand{\dd}{\,\mathrm d}\newcommand{\ii}{\mathrm i}
\newcommand{\R}{\mathbb R}\newcommand{\one}{\mathbf1}
\newcommand{\Sc}{\mathcal S}\newcommand{\norm}[1]{\lVert#1\rVert}
\begin{document}
{\color{ink}\LARGE\bfseries Die Schnittkopplungen rekonstruieren}\par
{\large Ausführbare Verklebung von Gravitation, Staub und Skalarfeld}\par
\textbf{Version 0.1, 11. November 2026.}\\
Projekt: \texttt{thomasmgfeller/x6-spectral-astrolabium}.\\
Quellstand: \texttt{efe19fca88a49d3d59de041652a60c26c180b8eb}.

\section*{Ergebnis}
Für die deklarierte endliche Kontrollfamilie liegt ein ausführbarer erweiterter
Randvertrag vor. Er behält die Schnittinzidenzen, passende Intervallzählungen
und die noch benötigten Uhr- und Feldvariablen als gemeinsame bedingte Daten.
Daraus rekonstruiert er die gemeinsame Wirkung exakt:
\[
 S_C=S_L+S_R-S_S+\Delta S.
\]
Das gemeinsame Einflussfunktional ergänzt die fünf zuvor nachgewiesenen Terme
und enthält außerdem die Maßkorrektur und einen einmaligen Schnittregulator:
\[
 \boxed{\ \mathcal I_\epsilon
 =2\exp\!\left[\frac{\ii}{\hbar}(\Delta S-S_S)\right]
       \exp\!\left[-\frac{\epsilon}{2}(q_3^2+q_4^2)\right].\ }
\]
Es multipliziert die \emph{bedingten lokalen Integranden}, bevor frühere
gekoppelte Variablen integriert werden. Bereits auf zwei Schnittpunkte
reduzierte Hilbert-Kerne enthalten die dafür nötigen Informationen nicht.

\begin{center}\begin{tabular}{p{.48\linewidth}p{.43\linewidth}}\toprule
Prüfung & Ergebnis\\\midrule
Gemeinsame Wirkung aus gespeicherten Daten & Exakte Übereinstimmung für alle $33^2=1089$ Paare.\\
Staub und Dichtebelegungen & Identität aller Wirkungskoeffizienten; gilt für alle Strata und Feldwerte.\\
Skalare Integrationsreihenfolge bei $\epsilon>0$ & 2178 zertifizierte Kontrollen bestanden.\\
Singuläre innere Skalarblöcke & 32 von 1089; als Distributionen in den Randfeldern definiert.\\
Reine Zweipunktinformation am Schnitt & Unzureichend; frühere Uhr- und Feldinformation ist nötig.\\
Beliebige globale CST-Historien & Nicht erfasst: zusätzliche Schnittumgehungen benötigen eigene Daten.\\\bottomrule
\end{tabular}\end{center}

\textbf{Interpretation.} Ein exakter endlicher Verklebungsvertrag ist erreicht.
Die bisher widerlegte Deutung als $L^2$-Evolution wird dadurch nicht wiederhergestellt.
Die neue gemeinsame Amplitude wird über Testzustandspaarungen definiert.
Weder eine physikalische vierdimensionale CST-Summe noch ein Bounce-Gesetz
wird damit konstruiert.

\textbf{Lieferumfang.} Python-Rekonstruktor ohne numerische Zusatzabhängigkeiten,
serialisierte Beispielpakete, unabhängige globale Kontrollrechnung, Arb-Zertifikate,
zehn Testgruppen, eingefrorene Quellen und Hashmanifest.
\newpage

\section{Unveränderte Wirkung und die fünf fehlenden Terme}
Wir verwenden die dimensionslosen Größen der bisherigen Kontrollen:
$\beta=4/\sqrt6$, $\kappa=1$, Skalarmasse null und ursprüngliche
Geometriekopplung eins. Für ein offenes Intervall $I(i,j)$ gilt
\[
 c(n)=
 \begin{cases}
 1&n=0,\\-9&n=1,\\16&n=2,\\-8&n=3,\\0&n\ge4,
 \end{cases}
 \qquad
 F_4(C)=|C|-\sum_{i\prec j}c(|I(i,j)|).
\]
Mit $c_{ij}=c(|I(i,j)|)$ lautet die vollständige Kontrollwirkung
\[
 \frac{S_C}{\hbar}
 =\kappa F_4(C)-\frac12\sum_{j\in C}r_j
 -\frac{\beta}{4}\sum_{i\prec j}c_{ij}
       \bigl[r_j(t_j-t_i)^2+(q_j-q_i)^2\bigr].                 \tag{1}
\]
Die Dichte sitzt am späteren Endpunkt. Der Skalarteil ist
$\beta q^TH^{(0)}q/2$; jedes Paar trägt
$H_{ii},H_{jj}\mathrel{-}=c_{ij}/2$ und
$H_{ij},H_{ji}\mathrel{+}=c_{ij}/2$ bei.

Die beiden Teilgebiete sind
\[
 L=\{0,1,2,3,4\},\quad R=\{3,4,5,6,7\},\quad S=L\cap R=\{3,4\}.
\]
Der Schnitt $S$ ist eine Antikette, in $L$ maximal und in $R$ minimal.
Die gemeinsame Ordnung ist ausdrücklich die \emph{minimale transitive Hülle}
der beiden lokalen Ordnungen. Zusätzliche Beziehungen, die nicht durch
$S$ vermittelt werden, sind in diesem Vertrag ausgeschlossen.

Für die vollständig geschichtete Ordnung mit Schichten
$(0,0,1,2,2,3,4,4)$ ergeben sich exakt:
\begin{center}\begin{tabular}{cccc}\toprule
Paar $(a,b)$ & $|I(a,b)|$ & $c_{ab}$ & Beitrag zu $\Delta F_4$\\\midrule
$(0,5)$ & 3 & $-8$ & $8$\\
$(1,5)$ & 3 & $-8$ & $8$\\
$(2,5)$ & 2 & $16$ & $-16$\\
$(2,6)$ & 3 & $-8$ & $8$\\
$(2,7)$ & 3 & $-8$ & $8$\\\bottomrule
\end{tabular}\end{center}
Die vier Paare $(0,6),(0,7),(1,6),(1,7)$ haben Intervallgröße vier
und Koeffizient null. Ihre Ordnungsrelation bleibt dennoch relevant.

Schreibt man $X$ für alle schnittübergreifenden Paare, ist
\[
 \Delta F_4=-\sum_{(a,b)\in X}c_{ab},\qquad
 \frac{\Delta S}{\hbar}
 =\kappa\Delta F_4-\frac{\beta}{4}
 \sum_{(a,b)\in X}c_{ab}
 \bigl[r_b(t_b-t_a)^2+(q_b-q_a)^2\bigr].                      \tag{2}
\]
Im konkreten Fall: $F_L=F_R=37$, $F_S=2$, $\Delta F_4=16$ und
$F_C=88$. Ohne die fünf Terme erhält man fälschlich $72$.
Die gemischte Ableitung
$\partial_{q_0}\partial_{q_5}(\Delta S/\hbar)=-4\beta$
zeigt direkt, weshalb ein Gegenterm nur in $q_3,q_4$ nicht genügt.
\newpage

\section{Rekonstruktion der Intervalle aus Randgedächtnis}
Setze $A=L\setminus S$ und $B=R\setminus S$. Jeder alte Endpunkt $a\in A$
speichert seinen Schnittfußabdruck, jeder neue Endpunkt $b\in B$ den
entsprechenden Eingangsfußabdruck:
\[
 U_a=\{s\in S:a\prec_L s\},\qquad
 V_b=\{s\in S:s\prec_R b\}.
\]
Zusätzlich werden für nichtleere Teilmengen $U,V\subseteq S$ gespeichert:
\begin{align*}
 H_L(a,U)&=\#\{z\in A:a\prec_L z,\ U_z=U\},\\
 H_R(b,V)&=\#\{z\in B:z\prec_R b,\ V_z=V\}.
\end{align*}
Jeder dieser Zähler darf separat bei vier abgeschnitten werden.
Die leere Teilmenge trägt niemals zu einem Kreuzintervall bei.

\textbf{Rekonstruktionssatz.}
Für die minimale transitive Verklebung gilt
\[
 a\prec_C b\quad\Longleftrightarrow\quad U_a\cap V_b\ne\varnothing .
\]
Wenn diese Bedingung erfüllt ist, berechnet man
\[
 \boxed{
 \begin{aligned}
 n_{ab}=\;&
 \sum_{U:\,U\cap V_b\ne\varnothing} H_L(a,U)
 +|U_a\cap V_b|\\
 &+\sum_{V:\,U_a\cap V\ne\varnothing} H_R(b,V).
 \end{aligned}}                                                 \tag{3}
\]
Für die Wirkung genügt $\widehat n_{ab}=\min(4,n_{ab})$ und
$c_{ab}=c(\widehat n_{ab})$.

\textbf{Beweis.}
Jeder Weg von $A$ nach $B$ muss den Schnitt durchlaufen. Wegen der
Antiketten- und Orientierungsbedingung kann er nicht wieder in $A$
zurückkehren. Dies beweist die Fußabdruckbedingung. Die Punkte eines offenen
Kreuzintervalls liegen disjunkt in $A$, $S$ oder $B$.
Ein $z\in A$ gehört genau dann dazu, wenn $a\prec_L z$ und
$U_z\cap V_b\ne\varnothing$; die erste Summe zählt genau diese Punkte.
Die Schnittpunkte sind exakt $U_a\cap V_b$. Die dritte Summe folgt analog
für $B$. Die lokalen Intervalle innerhalb $L$ und $R$ ändern sich nicht.
Da alle Zähler nichtnegativ sind,
$\min(4,\sum_j\min(4,h_j))=\min(4,\sum_j h_j)$.
Damit bewahrt die Abschneidung alle Wirkungskoeffizienten. \hfill$\square$

Im geschichteten Beispiel besitzen alle sechs Endpunkte den Fußabdruck
$S$. Die einzigen nichtverschwindenden Zähler sind
\[
 H_L(0,S)=H_L(1,S)=1,\qquad H_R(6,S)=H_R(7,S)=1.
\]
Alle übrigen sind null. Formel (3) ergibt etwa
$n_{05}=1+2+0=3$, $n_{25}=0+2+0=2$ und $n_{06}=1+2+1=4$.
So entstehen genau die fünf fehlenden Terme und die vier Nullkoeffizienten.

\textbf{Ausführbare Repräsentation.}
Die Schnittpunkte werden geordnet identifiziert. Teilmengen von $S=\{3,4\}$
werden als Bitmasken $1,2,3$ kodiert. \texttt{sewing.py} berechnet
Kreuzrelationen und Intervalle ausschließlich aus den Profilen.
Die Pakete enthalten zusätzlich die lokalen Relationstabellen samt
Wirkungskoeffizienten, einschließlich Nullkoeffizienten. Damit prüft die
Validierung jeden gespeicherten lokalen Zähler ohne globale Ordnungsabfrage.
\newpage

\section{Konkreter erweiterter Randvertrag}
Ein Paket ist weder eine neue Wellenfunktion auf zwei Punkten noch eine
gemittelte Geometrie. Es beschreibt eine bedingte lokale Amplitude mit explizit
gebundenen Gedächtnisvariablen. Historienlabels und Korrelationen bleiben erhalten.

\begin{center}\begin{tabular}{p{.29\linewidth}p{.62\linewidth}}\toprule
Gespeicherte Daten & Zweck und Bindung\\\midrule
Schnitt-IDs, Orientierung, äußere Rand-IDs & Eindeutige Identifikation beider Schnittpunkte; nur dort dürfen sich die Teilgebiete überlappen.\\
Lokale Wirkung und Historiengewicht & Unveränderte Gravitation, Staub- und Skalarphasen; kein vorheriges Mitteln der Hesse-Matrizen.\\
Fußabdrücke und Intervallhistogramme & Rekonstruktion aller Kreuzkoeffizienten durch (3).\\
Links: $t_a,q_a$ für $a\in A$ & Frühere Uhr- und Feldwerte für $(t_b-t_a)^2,(q_b-q_a)^2$.\\
Rechts: $t_b,r_b,q_b$ für $b\in B$ & Spätere Felder und die spätere Dichte im Kreuzterm.\\
Am Schnitt: $t_s,r_s,q_s$ & Gemeinsame Variablen, Schnittwirkung und Maß.\\
Maß, Kopplungswerte, Regulatorbesitz & Fixiert Quotienten, Dichtestrata und einmalige Dämpfung jedes inneren Feldes.\\
Abschlussregel und Integrität & Minimale transitive Hülle ohne Umgehungen; Schema und SHA-256 gegen unbemerkte Änderungen.\\\bottomrule
\end{tabular}\end{center}

\textbf{Wann darf eine alte Variable verschwinden?}
Die Kreuzwirkung benötigt keine frühere Dichte $r_a$. Sie darf bedingt auf
den erhaltenen Variablen integriert werden, sofern lokale Wirkung und
Randpräparation entsprechend mitgeführt werden. Für $t_a$ und $q_a$ ist
eine unabhängige Mittelung im Allgemeinen unzulässig. Insbesondere dürfen
$t_2,q_2$ vor dem Einfügen des Einflussfunktionals nicht verloren gehen.
Gleiches gilt rechts für die vom Kreuzterm benötigten Größen.

Für \emph{die eine vollständig geschichtete Ordnung} lässt sich der alte
Feldbedarf in den sechs Größen
\[
 t_0+t_1,\quad t_0^2+t_1^2,\quad t_2,\qquad
 q_0+q_1,\quad q_0^2+q_1^2,\quad q_2
\]
ausdrücken. Dabei müsste die gemeinsame bedingte Amplitude auf diesen Größen
korrekt mittransportiert werden; sechs getrennte Erwartungswerte reichen nicht.
Der ausführbare Vertrag verwendet deshalb die einfachere Speicherung pro
Endpunkt. Ein für beliebig viele Verklebungen konstantes oder minimales
Gedächtnis wird nicht behauptet.

\textbf{Fehlende globale Information ist sichtbar.}
Die lokalen Historien mit Codes $(8476,8592)$ besitzen eine minimale
Verklebung mit $F_4=14$. Fügt man global die Beziehung $1\prec5$ hinzu und
schließt transitiv ab, bleiben beide lokalen Einschränkungen unverändert,
aber $F_4=22$. Kein Verfahren kann diese Wahl aus denselben lokalen Paketen
ableiten. Für solche Modelle muss ein Vertrag zusätzliche globale
Kreuzrelationen und gegebenenfalls deren Gewichte als Eingabe verlangen.

Die Prüfsummen belegen Integrität, keine physikalische Richtigkeit.
Der aktuelle Rekonstruktor verwirft unbekannte Maß- oder Skalarparameter,
fehlende Feldbindungen, falsche Orientierung und inkonsistente Intervallprofile.
\newpage

\section{Einflussfunktional, Maß und gemeinsame Rechnung}
Die elementaren Maße lauten unverändert
\[
 \dd\lambda(t)=\tfrac12\one_{[-1,1]}\dd t,\qquad
 \dd\nu(r)=\tfrac12\delta_0(\dd r)+\tfrac12\one_{(0,1)}\dd r,
 \qquad \dd m_j=\dd\lambda(t_j)\dd\nu(r_j)\dd q_j.
\]
Ein zweipunktiger Rand besitzt
$\dd\mu_A=\tfrac12\prod_{j\in A}\dd m_j$;
ein innerer gelabelter Punkt besitzt $\dd m_j$.
Daher $\dd m_3\dd m_4=2\dd\mu_S$.
Es gibt keine weitere Summe über Schnittpermutationen:
die gemeinsamen IDs sind bereits fest identifiziert.

Mit $W_L=\exp(\ii S_L/\hbar)$, $W_R=\exp(\ii S_R/\hbar)$ gilt punktweise
\[
 W_LW_R\exp[\ii(\Delta S-S_S)/\hbar]=\exp(\ii S_C/\hbar).
                                                               \tag{4}
\]
Der Schnitt ist eine masslose Antikette, also
\[
 S_S/\hbar=2\kappa-\tfrac12(r_3+r_4).
\]
Der positive Regulator gehört einmal zu jedem gemeinsamen inneren Punkt
$I=\{2,3,4,5\}$. Lokal besitzen $L$ und $R$ bereits die Faktoren
$e^{-\epsilon q_2^2/2}$ bzw.\ $e^{-\epsilon q_5^2/2}$.
Für den konkreten Fünftermfall lautet der zusätzlich nötige Faktor somit
\[
 \mathcal I_\epsilon
 =2e^{-\epsilon(q_3^2+q_4^2)/2}
 \exp\!\left\{\ii\left[
 14\kappa+\tfrac12(r_3+r_4)
 -\tfrac\beta4\sum_{(a,b)\in X}c_{ab}
 \{r_b(t_b-t_a)^2+(q_b-q_a)^2\}\right]\right\}.                \tag{5}
\]
Die Zahl $14$ ist $16-2$: Kreuzgravitation minus doppelt gezählte Schnittgravitation.

Für beliebige Testpräparationen $f$ am Rand $\{0,1\}$ und $g$ am Rand
$\{6,7\}$ definieren wir für $\epsilon>0$
\[
 \mathcal B_\epsilon(g,f)=
 \frac1{1089}\sum_{\ell,r=1}^{33}
 \int \overline{g}\,f\,e^{\ii S_{C(\ell,r)}/\hbar}
 e^{-\epsilon\sum_{j\in I}q_j^2/2}
 \dd\mu_{01}\dd\mu_{67}\prod_{j\in I}\dd m_j.                 \tag{6}
\]
Ersetzt man im Integral den gemeinsamen Integranden durch
$W_{L,\epsilon}W_{R,\epsilon}\mathcal I_\epsilon$ und
$\prod_{I}\dd m_j$ durch $\dd m_2\,\dd\mu_S\,\dd m_5$,
bleibt (6) exakt gleich. Das ist die ausführbare Verklebungsregel.

\textbf{Beweis der gemeinsamen Rechnung.}
Die lokalen Intervalle bleiben gleich; die zusätzlichen Paare sind genau (3).
Vertexterme auf $S$ werden zweimal gezählt und einmal abgezogen.
Dies beweist (4) in allen drei Sektoren für beliebige reelle Feldwerte.
Der Faktor zwei korrigiert das Maß, die Regulatoren ergeben gemeinsam
die vier inneren Dämpfungen. Für Schwartz-Randfelder und integrable
Staubpräparationen ist bei $\epsilon>0$ absolute Konvergenz gegeben.
Fubini erlaubt daher jede Integrationsreihenfolge. Die endliche Summe
ändert daran nichts. Insbesondere umfasst die Identität alle $2^8$
Atom-/Positivdichtestrata einschließlich ihrer ursprünglichen Gewichte.
\hfill$\square$

Die Produktgewichte $1/33^2$ definieren hier eine explizite neue
Kontrollfamilie aus den bisherigen lokalen Familien. Sie sind kein
hergeleitetes physikalisches Maß auf sämtlichen achtpunktigen Ordnungen.
\newpage

\section{Nullmoden, zulässige Paarungen und Fehlerbegriff}
Die gemeinsame skalare Matrix $H^{(0)}$ hat für 1057 Paare einen invertierbaren
inneren Block $H_{II}$. Bei 32 Paaren gilt exakt
\[
 \operatorname{rank}H_{II}=3,\qquad
 \operatorname{rank}\begin{pmatrix}H_{II}\\H_{BI}\end{pmatrix}=4.
                                                               \tag{7}
\]
Jede innere Nullmode koppelt also an die äußeren Felder
$b=(q_0,q_1,q_6,q_7)$. Ein überall gewöhnlicher punktweiser Fresnelkern
darf für diese 32 Summanden nicht unterstellt werden.

\textbf{Distributioneller Abel-Grenzwert.}
Schreibe $M=\beta H^{(0)}$ und zerlege die inneren Variablen orthogonal
in $x$ auf den nichtentarteten Eigenraum und $y$ auf den Nullraum.
Seien $\Lambda$ der invertierbare Block, $J_x$ und $J_0$ die Kopplungen
dieser Räume an $b$. Dann ergibt die skalare Integration im Distributionssinn
\[
 \begin{split}
 G_0(b)={}&
 \frac{(2\pi)^{(4-k)/2}}{\sqrt{|\det\Lambda|}}\,
 e^{\ii\pi\,\operatorname{sgn}(\Lambda)/4}
 e^{\frac{\ii}{2}b^T(M_{BB}-J_x^T\Lambda^{-1}J_x)b}\\
 &{}\times (2\pi)^k\,\delta^{(k)}(J_0b).
 \end{split}                                                \tag{8}
\]
mit $k=\dim\ker H_{II}$; hier $k=0$ oder $1$.
Die Zeilen von $J_0$ sind bei $k=1$ unabhängig nach (7).
Die Delta-Distribution ist deshalb ein wohldefinierter Rückzug unter
einer surjektiven linearen Abbildung. Orthogonale Nullraumkoordinaten fixieren
die Normierung; ein unnormierter Nullvektor darf nicht unverändert
für den Delta-Vorfaktor eingesetzt werden.

Die nichtentarteten Gaußintegrale konvergieren zu den Fresnel-Faktoren.
Für die Nullrichtungen gilt
$\int e^{-\epsilon |y|^2/2+\ii y^TJ_0b}\dd y
\to(2\pi)^k\delta^{(k)}(J_0b)$ in $\Sc'(\R^4)$.
Die verbleibende quadratische Phase erhält den Schwartz-Testraum.
Damit existiert für jede der 1089 Historien der gepaarte Abel-Grenzwert.

Beispiel $(8592,8716)$: In der Reihenfolge $(2,3,4,5)$ ist
$(-1,0,0,1)^T$ eine Nullrichtung.
Ihr Produkt mit $H_{IB}$ ist $(0,15/2,-15/2,0)$.
Die entstehende Randbedingung lautet also $q_1=q_6$.
Sie ist eine verallgemeinerte Amplitude, keine gewöhnliche Dichtefunktion.

\textbf{Zustandsraum.}
Mit dem bisherigen kompakten Uhr-/Dichtemaß $\eta_A$ wählen wir
\[
 \Phi_A=
 \bigl(L^2(\eta_A)\otimes_{\rm alg}\Sc(\R^2)\bigr)_{\rm inv}.
\]
Die Invarianz vertauscht beide Punkte samt Uhr, Dichte und Feld gemeinsam.
Die Topologie ist die induzierte projektive Tensorprodukttopologie
aus $L^2$-Norm und Schwartz-Seminormen.
Dieser Testraum ist dicht im bisherigen Rand-Hilbertraum.
Die Amplitude ist eine kontinuierliche sesquilineare Paarung
$\Phi_{\rm out}\times\Phi_{\rm in}$, entsprechend $\Phi_{\rm in}\to\Phi_{\rm out}'$.
Die endliche Staubintegration ist beschränkt; ihre Phasen hängen nicht von den
Skalarfeldern ab. Die 32 Nullmoden erfordern daher keine neue Staubmaßannahme.

\textbf{Fehlerbegriff.}
Für feste Testzustände ist
$E_\epsilon(g,f)=|\mathcal B_\epsilon(g,f)-\mathcal B_0(g,f)|$.
Für direkte und rekonstruierte Rechnung ist die Differenz exakt null,
sowohl bei jedem $\epsilon>0$ als auch im gepaarten Grenzwert.
Eine uniforme Operatornormrate im früheren gewichteten Raum wird hier
nicht behauptet. Fehler durch ausgelassene globale Historien sind unbestimmt.
\newpage

\section{Zertifizierte Zahlenkontrollen und Grenzen}
Zusätzlich zur exakten Wirkungsidentität wird die vollständige Skalarintegration
für alle 1089 Paare mit Arb bei 192 Bit gerechnet.
Der numerische Kontrollsektor setzt \emph{alle acht Dichten auf ihre Atome}.
Die äußeren Präparationen sind die normierten bisherigen Atom-/Grundpakete
\[
 \Psi(q,r)=2\sqrt2\,\one_{\{r_1=r_2=0\}}\,
 \pi^{-1/2}e^{-(q_1^2+q_2^2)/2}.
\]
Dies wählt auch im Inneren nur den Atomanteil; es ist ausdrücklich
\emph{nicht} die numerische Integration der Positivdichtesektoren.

Die Randquotienten, Paketfaktoren und acht Atomgewichte ergeben gemeinsam
$(1/2)^2(2\sqrt2)^2(1/2)^8=1/128$. Mit
\[
 Q_\epsilon=\operatorname{diag}(1,1,\epsilon,\epsilon,\epsilon,
 \epsilon,1,1)-\ii\beta H^{(0)}
\]
ist der historische Beitrag
\[
 A_{\epsilon,C}=\frac{e^{\ii F_4(C)}}{128}\,
 \frac{16\pi^3}{\prod_{k=1}^{8}\sqrt{p_k(Q_\epsilon)}}.       \tag{9}
\]
Die $p_k$ sind Gauß-Eliminationspivots. Bei positivem Regulator liegt ihr
Realteil strikt rechts von null. Für den Grenzwert wird zuerst in der
Reihenfolge $(0,1,6,7,2,3,4,5)$ eliminiert; jeder Grenzpivot liegt
zertifiziert außerhalb von null und des Wurzelschnitts.
Das Produkt der fortgesetzten Einzelwurzeln bewahrt die Fresnelphase.
Eine unkontrollierte Hauptwurzel der Gesamtdeterminante wird nicht verwendet.

\begin{center}\begin{tabular}{p{.43\linewidth}rr}\toprule
Kontrolle & Realteil & Imaginärteil\\\midrule
Festes Paar $(25500,25500)$, gemeinsam
 & $-0.000076122522$ & $-0.000071394860$\\
Dasselbe, nur lokale Wirkungen
 & $ 0.000327584152$ & $-0.000729801282$\\
1089er-Summe, Abel-Grenzwert
 & $-0.000069112936$ & $ 0.000167149114$\\
1089er-Summe, $\epsilon=1$
 & $ 0.000002097094$ & $ 0.000052809063$\\
1089er-Summe, $\epsilon=1/7$
 & $-0.000047106955$ & $ 0.000135871965$\\\bottomrule
\end{tabular}\end{center}
Jede angezeigte Komponente hat absolute Rundungs- und Einschließungsunsicherheit
kleiner als $10^{-12}$. Die vollständigen Arb-Intervalle stehen in den JSON-Daten.
Die Regulatordifferenzen der 1089er-Atomkontrolle zum Abel-Wert betragen
zertifiziert ungefähr $1.3470158\,10^{-4}$ bzw.\ $3.8242951\,10^{-5}$.
Sie werden nicht als Fehlergrenzen für beliebige Zustände ausgegeben.

Für $\epsilon=1,1/7$ wird zusätzlich die Reihenfolge
$(2,5,3,4,0,1,6,7)$ geprüft: alle 2178 Ergebnisse überlappen ihre
unabhängig berechneten Arb-Intervalle. In diesem Schritt bleiben
die Kreuzkopplungen Teil der gemeinsamen Matrix.

\textbf{Was ist abgesichert?}
Die exakte Koeffizientenidentität erfasst auch sämtliche Positivdichtesektoren,
ohne hochdimensionale Quadratur zu benötigen. Die Zahlenkontrolle prüft
zusätzlich Phase, Maß und Integrationsreihenfolge in einem klar begrenzten
Sektor. Die fünf Einzelterme, Schnittphase und Maßfaktor besitzen
negative Kontrollen. Hundert größere künstliche Ordnungen prüfen ferner,
dass das Abschneiden der Zähler bei vier die Rekonstruktion erhält.

\textbf{Offen bleibt:} physikalische Historiengewichte, zusätzliche globale
Kreuzrelationen, ein kontrollierter Grenzübergang zu größeren Familien,
Fehler einer weiteren Gedächtnisreduktion und beliebig oft iterierbare
Komposition als physikalischer Evolutionsoperator.
\newpage

\section{Reproduktion und Quellenvertrag}
Die Datei \texttt{sewing.py} benötigt nur die Python-Standardbibliothek.
Die beigelegten JSON-Pakete enthalten tatsächlich alle für den Beispiel-Join
benötigten lokalen Daten. Die direkte globale Ordnung wird nur in der
unabhängigen Prüfrechnung verwendet.

\begin{verbatim}
python sewing.py examples/left.json examples/right.json \
  --fields examples/fields.json --output /tmp/joined.json
python analysis.py
python -m unittest discover -s . -p 'test_*.py' -v
python verify.py
python build_pdf.py
\end{verbatim}
Die Ausgabedatei enthält die rekonstruierte gemeinsame Wirkung,
alle neun Kreuzintervalle samt Nullkoeffizienten, den Maßfaktor,
die Regulatorzuständigkeit und bei Feldbindung die exakten rationalen
Koeffizienten der gemeinsamen Wirkung sowie der zusätzlichen Schnittphase.
Die Beispielpakete tragen je Gewicht $1/33$; der Join ergibt $1/1089$.
Mit \texttt{--epsilon 1/7} wird außerdem das Einflussfunktional an den
gebundenen Feldern ausgewertet; diese optionale Gleitkommaausgabe ist
kein Arb-Zertifikat. Frühere Dichten sind im Beispiel als lokale
Hilfsvariablen explizit gebunden.

\begin{center}\begin{tabular}{p{.39\linewidth}p{.52\linewidth}}\toprule
Datei & Inhalt\\\midrule
\texttt{contract.json} & Zustandsraum, Maße, Abschlussregel, Anspruch und Ausschlüsse.\\
\texttt{sewing.py} & Paketprüfung, Profilrekonstruktion, Wirkung und Schnittphase.\\
\texttt{analysis.py} & Globale Ordnungsreferenz, 1089er-Prüfung, Nullmoden und Arb-Kontrollen.\\
\texttt{examples/*.json} & Vollständige linke/rechte Pakete und rationale Feldwerte.\\
\texttt{results/verification.json} & Ergebnisübersicht und vollständige Zahlenintervalle.\\
\texttt{results/history\_pairs.json} & Exakte Kennzahlen jedes Historienpaars.\\
\texttt{results/singular\_pairs.json} & Alle 32 singulären Paare und ihre Nullrichtungen.\\
\texttt{test\_sewing.py} & Zehn Gruppen positiver und negativer Kontrollen.\\
\texttt{inputs/}, \texttt{provenance.json} & Sechs unveränderte Quellartefakte mit Git-Blob-Identitäten.\\
\texttt{MANIFEST.sha256.json} & Vollständige Dateiintegrität des Prüfpakets.\\\bottomrule
\end{tabular}\end{center}

\textbf{Anschluss an den bisherigen roten Faden.}
Die Quelle \texttt{reduction\_gluing\_v0\_1} liefert die fünf Intervallterme
und das frühere feste Achtpunkt-Gegenbeispiel zur lokalen Verklebung.
\texttt{coupled\_kernel\_domain\_v0\_1} klärt den Definitionsbereich der
ursprünglichen 33er-Amplitude. Der vorliegende Vertrag ergänzt jetzt
die fehlende Gedächtnisstruktur und rekonstruiert die gemeinsame Wirkung.

Die eingefrorenen Dateien sind Referenzdaten; insbesondere ist
\texttt{inputs/prior\_gluing.py} ohne sein ursprüngliches Nachbarpaket
kein eigenständiges Programm. Die ausführbaren Einstiege dieses Pakets
sind oben aufgeführt. \texttt{verify.py} kontrolliert alle sechs
Git-Blob-Identitäten, Dateihashes, Tests und eine byteidentische
Neuberechnung der JSON-Ergebnisse in einem temporären Verzeichnis.

\textbf{Statusentscheidung.}
\texttt{FINITE\_EXTENDED\_SEWING\_PASS}: Der endliche erweiterte Randvertrag
reproduziert die gemeinsame Rechnung in der deklarierten Interpretation.\\
\texttt{PHYSICAL\_CST\_EVOLUTION\_NOT\_CONSTRUCTED}: Physikalische
Evolution, Unitarität, allgemeine Junction-Dynamik und Singularitätsauflösung
bleiben unbewiesen. Bestehende physikalische Gates werden nicht freigegeben.

\textbf{Lizenz.} Apache-2.0; unveränderte Quellartefakte gemäß dem
angegebenen Projektstand. Keine fremden Messdaten und keine kostenpflichtigen
Rechenläufe wurden benötigt.
\end{document}
